% 復習1　AIのハルシネーションを確かめる
\session{1}{1}{40分}{AIのハルシネーションを確かめる}{出力の検証}{}

\goals{
  \item この授業でのAI利用の基本（大学の Copilot を使う、出力を検証する、使ったことを記録する）を説明できる
  \item AIがもっともらしい誤り（ハルシネーション）を出す場面を、自分で確かめられる
  \item 出力を使う前に何を確かめるかを、自分の言葉で挙げられる
}

\work{確認}{AI利用ガイドの「基本の考え方」（自分の言葉で4つ）}
\writelines{4}

\work[12mm]{準備}{Copilot の動作確認}
\begin{checks}
  \item 大学のアカウントで Copilot にサインインできた
  \item 簡単な質問をして、回答が返ってくることを確かめた
\end{checks}

\work[70mm]{ワーク1}{AIに文献と数値を挙げさせ、確かめる}
\begin{promptbox}[ハルシネーションを体験する]
日本の中小企業における生成AIの導入状況について、参考文献を3つと具体的な数値を挙げて、300字で説明してください。
\end{promptbox}
\instr{AIが挙げた文献と数値が実在するか、CiNii Research、図書館の蔵書検索、官公庁の公式サイトで確かめる。}
\begin{wstabh}{colspec={X[3,l,m] Q[c,wd=14mm,m] X[2,l,m]},row{2-Z}={ht=10mm}}
  AIが挙げた文献・数値 & 実在 ○× & 確かめた場所・わかったこと \\
  & & \\ & & \\ & & \\ & & \\
\end{wstabh}

\work[50mm]{ワーク2}{自分がよく知っている話題で正誤を判定する}
\instr{アルバイト先の業界、出身地、部活、趣味など、自分が詳しく知っている話題について聞き、内容の正誤を判定する。}
\begin{wstabh}{colspec={X[2,l,m] X[3,l,m] Q[c,wd=14mm,m] X[2,l,m]},row{2-Z}={ht=12mm}}
  尋ねたこと & AIの答え（要点） & 正誤 ○△× & どこが違ったか \\
  & & & \\ & & & \\ & & & \\
\end{wstabh}

\work{まとめ}{気づいたこと・「出力を使う前のチェック」}
\instr{AI利用ガイドの「出力を使う前のチェック」も見ながら、自分の言葉で整理する。}
\writelines{5}

\begin{nexttime}
  \item 関心のあるテーマ（経営・情報・企業・働き方など）を3つ書き出す（復習2で使う）。「なぜ気になるのか」を一言添える。
\end{nexttime}
\begin{wstabh}{colspec={Q[c,wd=8mm,m] X[2,l,m] X[3,l,m]},row{2-Z}={ht=11mm}}
  & 関心のあるテーマ & なぜ気になるのか \\
  1 & & \\ 2 & & \\ 3 & & \\
\end{wstabh}
\aimemo
